\section{Penjumlahan dan Pengurangan}
\label{sec:penjumlahan-pengurangan}

Operasi aritmatika bekerja pada operand yang sudah berada di register
atau memori. Instruksi penjumlahan dan pengurangan menyimpan hasil
pada operand tujuan, sehingga operand tujuan berubah sedangkan
operand sumber tetap.

Bagian ini diringkas dari \cite{gunawan2019}, Bagian 14.2--14.3,
hlm.~132--137, dan \cite{sto2001}, BAB VIII, hlm.~32--40.

\subsection{ADD, ADC, dan INC}

\begin{description}
    \item[\asm{ADD}] Menjumlahkan operand sumber dengan operand
          tujuan. Hasil disimpan pada operand tujuan.
          \begin{center}
          \lstasm|ADD tujuan, sumber|
          \end{center}

    \item[\asm{ADC}] Menjumlahkan operand sumber, operand tujuan,
          dan nilai \flag{CF} dari operasi sebelumnya.
          \begin{center}
          \lstasm|ADC tujuan, sumber|
          \end{center}

    \item[\asm{INC}] Menambah operand tujuan dengan satu.
          \begin{center}
          \lstasm|INC tujuan|
          \end{center}
\end{description}

Contoh:

\begin{lstlisting}[language=MASM, numbers=none]
MOV AX, 25
ADD AX, 17          ; AX = 42
ADC AX, BX          ; AX = AX + BX + CF
INC AX              ; AX = AX + 1
\end{lstlisting}

\asm{ADC} dipakai ketika hasil penjumlahan sebelumnya menghasilkan
\eng{carry} yang harus dibawa ke word berikutnya. Pada penjumlahan
satu word, \asm{ADC} sering tampak sama dengan \asm{ADD} karena
\flag{CF} biasanya nol.

\subsection{SUB, SBB, dan DEC}

\begin{description}
    \item[\asm{SUB}] Mengurangi operand tujuan dengan operand sumber.
          \begin{center}
          \lstasm|SUB tujuan, sumber|
          \end{center}

    \item[\asm{SBB}] Mengurangi operand tujuan dengan operand sumber
          dan nilai \flag{CF} (pinjaman) dari operasi sebelumnya.
          \begin{center}
          \lstasm|SBB tujuan, sumber|
          \end{center}

    \item[\asm{DEC}] Mengurangi operand tujuan dengan satu.
          \begin{center}
          \lstasm|DEC tujuan|
          \end{center}
\end{description}

Contoh:

\begin{lstlisting}[language=MASM, numbers=none]
MOV AX, 100
SUB AX, 45          ; AX = 55
SBB AX, BX          ; AX = AX - BX - CF
DEC AX              ; AX = AX - 1
\end{lstlisting}

\subsection{Pengaruh terhadap Flag}

Setiap operasi aritmatika memperbarui \eng{flag}. Arti \eng{flag}
yang relevan dirangkum pada Tabel~\ref{tab:flag-aritmatika}.

\begin{table}[H]
    \centering
    \caption{Makna flag pada operasi aritmatika.}
    \label{tab:flag-aritmatika}
    \small
    \begin{tabular}{@{}p{1.4cm}p{7.2cm}p{5.2cm}@{}}
        \toprule
        \textbf{Flag} & \textbf{Arti} & \textbf{Tanda pada hasil} \\
        \midrule
        \flag{CF} & Ada simpanan keluar dari bit tertinggi atau
        pinjaman pada pengurangan &
        Hasil tak bertanda melampaui kapasitas \\
        \addlinespace
        \flag{OF} & Hasil bertanda di luar jangkauan &
        Luapan bilangan bertanda \\
        \addlinespace
        \flag{ZF} & Hasil operasi bernilai nol &
        Kedua operand sama besar \\
        \addlinespace
        \flag{SF} & Bit tanda hasil bernilai 1 &
        Hasil bertanda negatif \\
        \addlinespace
        \flag{AF} & Simpanan dari empat bit rendah &
        Dipakai pada operasi BCD \\
        \addlinespace
        \flag{PF} & Jumlah bit 1 pada byte rendah genap &
        Pemeriksaan paritas \\
        \bottomrule
    \end{tabular}
\end{table}

\flag{CF} dan \flag{OF} menjawab pertanyaan yang berbeda. \flag{CF}
dipakai untuk bilangan tidak bertanda, sedangkan \flag{OF} dipakai
untuk bilangan bertanda. Contoh \hex{FF} + \hex{01} menghasilkan
\hex{00}: \flag{CF} bernilai 1 karena ada simpanan keluar, tetapi
\flag{OF} bernilai 0 karena hasil bertanda $-1 + 1 = 0$ masih berada
di dalam jangkauan.

\subsection{Aritmatika Lebih dari 16 Bit}

Bilangan 32 bit dipecah menjadi dua word: word rendah (bit 0--15)
dan word tinggi (bit 16--31). Penjumlahan dimulai dari word rendah
dengan \asm{ADD}. \eng{Carry} yang dihasilkan dibawa ke word tinggi
melalui \asm{ADC}. Pola yang sama berlaku untuk pengurangan dengan
\asm{SUB} dan \asm{SBB}.

\begin{figure}[H]
    \centering
    \begin{tikzpicture}[
        kotak/.style={
            draw=obeBlue, line width=0.8pt, fill=obeBlueLight,
            rounded corners=3pt, minimum width=3.0cm,
            minimum height=0.9cm, align=center, font=\small,
            inner sep=3pt
        },
        hasil/.style={
            draw=obeAmber, line width=0.8pt, fill=obeAmberLight,
            rounded corners=3pt, minimum width=3.0cm,
            minimum height=0.9cm, align=center, font=\small,
            inner sep=3pt
        },
        panah/.style={-{Stealth[length=2.2mm]}, line width=0.8pt,
            obeTeal},
        ket/.style={font=\scriptsize\itshape, text=obeGray,
            align=center}
    ]
        \node[kotak] (alo) at (0, 2.4) {word rendah A};
        \node[kotak] (blo) at (0, 1.3) {word rendah B};
        \node[hasil] (rlo) at (5.2, 1.85) {hasil rendah};
        \draw[panah] (alo.east) -- (rlo.west);
        \draw[panah] (blo.east) -- (rlo.west);
        \node[ket] at (7.5, 1.85) {\asm{ADD}};

        \node[kotak] (ahi) at (0, -0.2) {word tinggi A};
        \node[kotak] (bhi) at (0, -1.3) {word tinggi B};
        \node[hasil] (rhi) at (5.2, -0.75) {hasil tinggi};
        \draw[panah] (ahi.east) -- (rhi.west);
        \draw[panah] (bhi.east) -- (rhi.west);
        \node[ket] at (7.5, -0.75) {\asm{ADC}};

        \draw[panah, obeMaroon, line width=1pt]
            (rlo.south) -- (rhi.north)
            node[midway, left, font=\scriptsize, text=obeMaroon]
            {\flag{CF}};
    \end{tikzpicture}
    \caption{Propagasi \eng{carry} dari word rendah ke word tinggi
    pada penjumlahan 32 bit.}
    \label{fig:carry-32bit}
\end{figure}

Contoh penjumlahan 32 bit:

\begin{lstlisting}[language=MASM, numbers=none]
MOV AX, rendahA
ADD AX, rendahB     ; hasil rendah, CF mungkin 1
MOV hasilRendah, AX
MOV AX, tinggiA
ADC AX, tinggiB     ; carry dari langkah pertama ikut dijumlahkan
MOV hasilTinggi, AX
\end{lstlisting}

Bila tahap kedua memakai \asm{ADD}, nilai \flag{CF} dari tahap
pertama hilang dan hasilnya berkurang satu.

\begin{peringatanbox}
\vspace{1.1em}
\asm{INC} dan \asm{DEC} tidak mengubah \flag{CF}. Karena itu,
keduanya tidak boleh dipakai untuk meneruskan \eng{carry} pada
aritmatika bertingkat. Gunakan \asm{ADD} atau \asm{ADC} bila
\flag{CF} perlu diperbarui.
\end{peringatanbox}
